% Two disjoint intra-communicators (say, two libraries) bridged by an
% inter-communicator, which is what lets a rank in one address the other.
\documentclass[dvisvgm,border=3pt]{standalone}
\input{preamble.tex}
\begin{document}
\begin{tikzpicture}[x=1cm,y=1cm,
    rk/.style   ={circle,fill=uibknavy,draw=blu,line width=0.5pt,text=onfill,
                  font=\sffamily\scriptsize,minimum size=5.6mm,inner sep=0pt},
    inter/.style={ellipse,fill=blu!45,draw=blu,line width=0.6pt},
    intra/.style={ellipse,fill=uibkorange,draw=uibkorange!70!black,line width=0.6pt},
    note/.style  ={text=fg,font=\sffamily\scriptsize}]

  % the inter-communicator spans both sides
  \node[inter,minimum width=46mm,minimum height=20mm] at (-1.55,0) {};
  \node[inter,minimum width=46mm,minimum height=20mm] at (1.55,0) {};

  \foreach \s in {-1,1}{
    \node[intra,minimum width=34mm,minimum height=14mm] at ({\s*1.55},0) {};}

  \foreach \i/\x/\y in {0/-0.72/-0.22, 1/-0.06/0.26, 2/0.44/-0.30, 3/1.06/0.16}{
    \node[rk] at ({\x-1.55},\y) {\i};
    \node[rk] at ({\x+1.55},\y) {\i};}

  % the boundary between the two libraries
  \draw[draw=fg,line width=1.0pt,dash pattern=on 4pt off 3pt] (0,1.30) -- (0,-1.30);
  \node[note,anchor=south east] at (-0.08,1.32) {library A};
  \node[note,anchor=south west] at (0.08,1.32) {library B};

  \draw[->,draw=blu,line width=0.9pt] (-3.55,1.05) -- (-2.30,0.34);
  \node[note,anchor=south east] at (-3.50,1.02) {intra-communicator rank};
  \draw[->,draw=blu,line width=0.9pt] (3.70,0.95) -- (2.95,0.52);
  \node[note,anchor=west] at (3.74,0.98) {intra-communicator};
  \draw[->,draw=blu,line width=0.9pt] (3.70,-1.05) -- (2.80,-0.72);
  \node[note,anchor=west] at (3.74,-1.08) {inter-communicator};

\end{tikzpicture}
\end{document}
